\question{4.8}{
    Een verzekeringsmaatschappij heeft een analyse gemaakt van de schadeclaims die in een gegeven jaar werden ingediend door mensen met een WA-verzekering.
    Op basis van deze gegevens is de kansverdeling geformuleerd in de tabel.
    Hierin staat per individuele verzekerde aangegeven hoe groot in een gegeven jaar de kansen zijn op uitbetaling van een bepaald schadebedrag.
    \begin{center}
        \begin{tabular}{cc}
            \toprule
                {\bfseries Schadebedrag $k$ (euro)} & {\bfseries Kans $P(X=k)$} \\
            \cmidrule{1-1} \cmidrule{2-2}
                $0$ & $0.85$ \\
                $100$ & $0.05$ \\
                $200$ & $0.04$ \\
                $400$ & $0.03$ \\
                $1000$ & $0.02$ \\
                $3000$ & $0.01$ \\
            \bottomrule
        \end{tabular}
    \end{center}
}

    \begin{enumerate}[label=(\alph*)]
        \item Bereken het verwachte schadebedrag per klant.
        \answer{
            Het verwachte schadebedrag per klant is de verwachtingswaarde van de kansvariabele $X$ en wordt als volgt berekend:
            \begin{align*}
                E(X)    &= \sum_k k \cdot P(X=k) \\
                        &= 0 \cdot P(X=0) + 100 \cdot P(X=100) + \ldots + 3000 \cdot P(X=3000) \\
                        &= 0 \cdot 0.85 + 100 \cdot 0.05 + \ldots + 3000 \cdot 0.01 \\
                        &= 75.
            \end{align*}
            Het verwachte schadebedrag per klant is dus \euro $75$.
        }

        \item Veronderstel dat de schadeclaims per klant voor opeenvolgende jaren mogen worden beschouwd als onderling onafhankelijke kansvariabelen.
        Hoe groot is de kans dat een willekeurige klant in drie opeenvolgende jaren telkens een bedrag van minstens $1000$ euro claimt?
        Hoe groot is de kans dat men in drie opeenvolgende jaren geen enkele keer een schadebedrag claimt (dus $X=0$ in alle jaren)?
        \answer{
            De kans dat een willekeurige klant in drie opeenvolgende jaren telkens een bedrag van minstens \euro $1000$ claimt, is de kans dat $X$ in elk van de drie jaren gelijk is aan $1000$ of $3000$.
            We berekenen deze kans (voor een enkel jaar) als volgt:
            \begin{align*}
                P(X \ge 1000)   &= P(X=1000) + P(X=3000) \\
                                &= 0.02 + 0.01 \\
                                &= 0.03.
            \end{align*}
            De kans dat een klant in drie opeenvolgende jaren telkens een bedrag van minstens \euro $1000$ claimt, is dan (aangezien de claims per jaar onafhankelijk zijn):
            \[
                P(X \ge 1000)^3 = (0.03)^3 = 0.000027.
            \]

            Op eenzelfde manier is de kans dat een klant in een enkel jaar geen schadebedrag claimt $P(X=0) = 0.85$.
            Aangezien de claims per jaar onafhankelijk zijn, is de kans dat een klant in drie opeenvolgende jaren geen enkel schadebedrag claimt:
            \[
                P(X=0)^3 = (0.85)^3 \approx 0.6141.
            \]
            De kans dat men in drie opeenvolgende jaren geen enkele keer een schadebedrag claimt is dus ongeveer $61.41\%$.
        }

        \item De WA-verzekering kost \euro $98$ per verzekerde per jaar.
        De maatschappij heeft $25000$ klanten.
        Hoe is de verwachte winstbijdrage per jaar?
        \answer{
            We kunnen de verwachtewinstbijdrage per klant berekenen als het verschil tussen de premie en het verwachte schadebedrag per klant (berekend in deel (a)):
            \[
                \text{Winstbijdrage per klant} = \text{Premie} - E(X) = 98 - 75 = 23.
            \]
            Aangezien de maatschappij $25000$ klanten heeft, is de verwachte winstbijdrage per jaar:
            \[
                \text{Totale winstbijdrage} = \text{Aantal klanten} \cdot \text{Winstbijdrage per klant} = 25000 \cdot 23 = 575000.
            \]
        }
    \end{enumerate}